\section{Exact Exponential Solution}
\label{sec:exact-solution}

The paraxial construction identifies the exponential profile
\begin{equation}
    n(h)=n_0 e^{\kappa h}
    \label{eq:exact-exponential-profile}
\end{equation}
as the natural candidate for reproducing the spherical horizon law above a
flat boundary. The first integral derived in Sec.~\ref{sec:model} also makes
this profile exactly solvable, without invoking the small-slope approximation.

\subsection{Exact ray trajectory}

For a ray whose minimum height is $h_m$, the slope vanishes at
$x=x_m$. Evaluating the first integral there gives
\begin{equation}
    a=n(h_m)=n_0 e^{\kappa h_m}.
\end{equation}
Substitution into the exact slope equation yields
\begin{equation}
    \left(\frac{dh}{dx}\right)^2
    =
    \frac{n^2(h)}{n^2(h_m)}-1
    =
    e^{2\kappa(h-h_m)}-1.
    \label{eq:exact-exponential-slope}
\end{equation}
On either branch of the ray,
\begin{equation}
    \left|x-x_m\right|
    =
    \int_{h_m}^{h}
    \frac{d\eta}
    {\sqrt{e^{2\kappa(\eta-h_m)}-1}}.
\end{equation}
Evaluating the integral gives
\begin{equation}
    \left|x-x_m\right|
    =
    \frac{1}{\kappa}
    \arccos\!\left[e^{-\kappa(h-h_m)}\right].
    \label{eq:exact-horizontal-distance}
\end{equation}
Inverting Eq.~\eqref{eq:exact-horizontal-distance} produces the exact ray
family
\begin{equation}
    h(x)
    =
    h_m
    -
    \frac{1}{\kappa}
    \ln\!\left[
        \cos\!\bigl(\kappa(x-x_m)\bigr)
    \right],
    \label{eq:exact-exponential-ray}
\end{equation}
valid on the interval
\begin{equation}
    \left|\kappa(x-x_m)\right|<\frac{\pi}{2}.
\end{equation}
Near the minimum, expansion of the logarithm gives
\begin{equation}
    h(x)
    =
    h_m
    +
    \frac{\kappa}{2}(x-x_m)^2
    +
    \frac{\kappa^3}{12}(x-x_m)^4
    +
    O\!\left(\kappa^5(x-x_m)^6\right),
    \label{eq:exact-ray-local-expansion}
\end{equation}
which recovers the parabolic result of the previous section at leading
order. The exact curvature is not constant away from the minimum. Indeed,
Eq.~\eqref{eq:exact-exponential-ray} implies
\begin{equation}
    h''
    =
    \kappa\sec^2\!\bigl(\kappa(x-x_m)\bigr)
    =
    \kappa\left(1+h'^2\right),
    \label{eq:exact-ray-curvature-equation}
\end{equation}
which reduces to $h''\simeq\kappa$ only in the paraxial regime.

\subsection{Exact grazing distance}

A limiting ray that just touches the opaque boundary has
\begin{equation}
    h_m=0.
\end{equation}
The horizontal distance from the tangency point to an endpoint at height $h$
is therefore
\begin{equation}
    d_{\mathrm{e}}(h)
    =
    \frac{1}{\kappa}
    \arccos\!\left(e^{-\kappa h}\right).
    \label{eq:exact-exponential-horizon-distance}
\end{equation}
For observer and target heights $h_{\mathrm{o}}$ and $h_{\mathrm{t}}$, the
exact limiting separation is
\begin{equation}
    L_{\max}^{(\mathrm{e})}
    =
    \frac{1}{\kappa}
    \left[
        \arccos\!\left(e^{-\kappa h_{\mathrm{o}}}\right)
        +
        \arccos\!\left(e^{-\kappa h_{\mathrm{t}}}\right)
    \right].
    \label{eq:exact-exponential-maximum-separation}
\end{equation}
With $\kappa=1/R$, this becomes
\begin{equation}
    L_{\max}^{(\mathrm{e})}
    =
    R
    \left[
        \arccos\!\left(e^{-h_{\mathrm{o}}/R}\right)
        +
        \arccos\!\left(e^{-h_{\mathrm{t}}/R}\right)
    \right].
    \label{eq:exact-earth-scale-separation}
\end{equation}
Thus the exponential profile produces a genuine boundary-occlusion threshold
above a flat surface: beyond the limiting separation, every connecting ray
intersects
the boundary, whereas raising either endpoint restores a non-intersecting
ray.

\subsection{Paraxial limit}

For $z=\kappa h\ll 1$,
\begin{equation}
    \arccos\!\left(e^{-z}\right)
    =
    \sqrt{2z}
    \left(
        1
        -\frac{z}{6}
        +\frac{z^2}{120}
        +O(z^3)
    \right).
    \label{eq:arccos-exponential-expansion}
\end{equation}
Consequently,
\begin{equation}
    d_{\mathrm{e}}(h)
    =
    \sqrt{\frac{2h}{\kappa}}
    \left(
        1
        -\frac{\kappa h}{6}
        +\frac{(\kappa h)^2}{120}
        +O\!\left((\kappa h)^3\right)
    \right).
    \label{eq:exact-distance-expansion}
\end{equation}
The leading term is precisely the paraxial distance in
Eq.~\eqref{eq:paraxial-horizon-distance}. The higher-order terms quantify
where the exponential refractive construction departs from a constant-
curvature approximation. The next section compares this exact optical law
with the corresponding horizon law on a spherical surface.
